\section{Introduction}

Formula 1 is the pinnacle of motorsport, combining engineering excellence with
driver skill in a high-stakes global competition. Each race involves 20 drivers
from 10 constructor teams competing across circuits with varying characteristics,
weather conditions, and strategic possibilities. The outcome of any given race
depends on a complex interaction of factors: qualifying performance, race pace,
tire strategy, pit stop timing, mechanical reliability, and driver decision-making
under pressure.

Predicting race outcomes has significant value for multiple stakeholders: teams
use predictive models to optimize strategy, broadcasters enhance viewer experience
with real-time probability estimates, and analysts study the relative importance
of car versus driver performance. Despite this value, public research on F1
outcome prediction remains limited compared to other sports.

\subsection{Research Questions}

This paper addresses three primary research questions:

\begin{enumerate}
    \item Can machine learning models significantly outperform a grid-position
          baseline for predicting F1 race finishing positions?
    \item Which feature categories (lap telemetry, driver history, constructor
          form, weather) contribute most to prediction accuracy?
    \item Does temporal validation reveal systematic degradation in model
          performance as the season progresses?
\end{enumerate}

\subsection{Contributions}

Our primary contributions are:
\begin{itemize}
    \item A complete, reproducible ML pipeline using the FastF1 API for
          automated data extraction.
    \item A multi-category feature engineering framework with 50+ domain-specific
          features.
    \item Temporal cross-validation methodology that respects the sequential nature
          of racing seasons.
    \item Systematic comparison of gradient boosting, deep learning, and baseline
          approaches on the same dataset and evaluation protocol.
\end{itemize}

\subsection{Paper Structure}

Section~\ref{sec:methodology} describes data sources, feature engineering, model
architectures, and evaluation protocols. Section~\ref{sec:results} presents
quantitative results and feature importance analysis. Section~\ref{sec:conclusions}
summarizes findings and discusses directions for future work.
